% §1 Introduction
Conversational voice agents are now operational infrastructure in banking.
Production deployments handle account inquiries, card operations, transfers, and
authentication over the telephone, at cadences measured in millions of calls per year \citep{nubank2026}.
The institutions that operate them are not judged on how gracefully a single turn is handled.
They are judged on conversation-level key performance indicators:
the containment rate (calls resolved without human involvement),
the abandonment rate (callers who hang up mid-interaction),
the handoff rate (transfers to human agents),
and the compliance-escape rate (unauthorized commitments, disclosure omissions, undetected fraud).
Every one of these indicators is a property of the \emph{whole conversation},
yet the research machinery we hold for voice agents is almost entirely \emph{turn-level}.

This mismatch between what is measured and what is modeled is the starting point of this paper.
The voice-agent literature formalizes single decisions:
when the user has finished speaking \citep{tauvoice2026,x2turn2026},
how fast a response must begin \citep{garcia2026latencyfloor},
how quickly an agent must stop talking when interrupted \citep{garcia2026cessationfloor}.
Agent benchmarks score tool-using dialogues by full-episode success---whether the final
environment state is correct---treating the trajectory itself as data rather than as the object
of analysis \citep{sierra2024taubench}.
Workflow-extraction methods recover the structure of completed service conversations,
but structure without outcomes \citep{salesforce2025workflows}.
Escalation policy is decided from queue-level state, outside the conversation \citep{queuecontrol2026}.
No formal model connects the partial transcript of an in-flight voice conversation
to the terminal outcome that the business will record when the call ends.

Meanwhile, the failures that matter in production emerge progressively, not abruptly.
Recent work on task-oriented dialogue shows that breakdowns accumulate through
slot instability, unresolved contradictions, and task drift long before the final
outcome makes them visible \citep{nourbakhsh2026tracer}.
In banking this progression has names the industry tracks: the caller who repeats
an account number three times, the tool that fails twice on the same operation,
the caller who goes silent while the agent keeps narrating options.
Each is a trajectory pattern, and each becomes actionable only if it is detected
\emph{while the conversation can still be changed}.
The nearest neighbor to this problem, TRACER, forecasts eventual dialogue failure from
partial context with a learned classifier and a thresholded recovery trigger \citep{nourbakhsh2026tracer}.
It establishes that partial dialogues carry outcome signal from as little as 25\% of the
conversation.
But its trigger is a dev-set-tuned threshold with no decision-time guarantee,
its outcome is a single binary failure flag,
its evidence base is English text-chat benchmarks,
and its pipeline contains none of the physics that makes voice conversations hard:
endpointing errors, barge-ins, tool round-trips, authentication events.

The gap has two halves, corresponding to gaps CD1 and CD2 of our research program.
First, no trajectory grammar exists for voice-agent conversations:
a formal state space over conversation states (authentication, intent, tool loop,
clarification, confirmation) with typed transitions,
terminal outcome classes aligned with business KPIs
(contained, abandoned, handoff, compliance-escape),
and a drift statistic that formally distinguishes a conversation that is merely
\emph{longer} from one that is \emph{off-path}.
Second, no decision rule exists with a guarantee:
a stopping rule that, given the posterior over outcome classes induced by the partial
transcript, commits to an outcome prediction (and an intervention) at a calibrated
confidence within a bounded expected number of turns.
Sequential analysis provides the right machinery---sequential probability ratio
tests stop LLM reasoning paths early with calibrated error control \citep{consol2025}---but
it has never been transposed from reasoning paths inside one model call
to multi-turn conversations between an agent and a customer.

This paper closes both gaps with machinery our program has already published,
extended from the turn scale to the conversation scale.
We previously proved that cascade voice pipelines have an irreducible latency floor
$\Lambda^*$ and that tool-calling agents compound it through $k$-step reasoning loops \citep{garcia2026latencyfloor}.
We proved that response \emph{cessation}---how fast an agent can stop talking when a
barge-in occurs---has its own floor, structured as detect/decide/cease/resume \citep{garcia2026cessationfloor}.
We showed that text-only turn-end detection in Spanish is near chance on hard cuts,
which forces any conversation-level state estimator to condition on intent and tool
state rather than on the transcript alone \citep{garcia2026asking}.
These three results are load-bearing here.
The $k$-step ReAct bound becomes the cycle-length constraint of the grammar:
a run that exceeds $k^*$ tool-loop cycles is off-path with high probability.
The latency floor and cessation floor convert ``early enough to act'' from a slogan
into an accounting identity: an intervention trigger fires through the same pipeline
whose components have measured latencies, and the intervention itself must clear
the cessation floor to take the floor from the agent mid-utterance.
The Spanish partial-transcript pipeline supplies the sensing substrate and warns
against over-trusting text.

We make the following contributions.

\begin{enumerate}
\item \textbf{Trajectory grammar} (\S\ref{sec:grammar}).
We define a formal state space $G = (Q, \Sigma, \delta)$ over voice-agent
conversations, with seven interior states (greeting, auth, intent, tool loop,
clarification, confirmation, close) and four terminal outcome classes
(contained, abandoned, handoff, compliance-escape).
Transitions are typed (user-answer, user-bargein, agent-clarify, tool-call,
tool-fail, endpoint-error, disclosure, escalation), so that voice physics and
tool events are first-class grammar emissions.
Regressions are grammar cycles, and we define cycle depth $D(\sigma)$ as the
drift statistic separating long from off-path conversations,
bounded by the published $k$-step ReAct constraint \citep{garcia2026latencyfloor}.

\item \textbf{Trajectory Stopping Theorem} (\S\ref{sec:stopping}).
We cast early outcome prediction as partial-observation sequential decision-making.
Given the posterior $p(\theta \mid x_{1:k})$ over outcome classes
$\theta \in \mathcal{T}$ induced by the partial transcript and tool-event log,
a confidence-sequence stopping rule achieves expected decision time
$\mathbb{E}[k_{dec}] \leq k^*(\delta, \text{costs})$ at confidence $1-\delta$.
We derive false-early-exit economics from a banking cost model in which false
containment costs more than false handoff and compliance escape has unbounded cost,
and we show that TRACER's fixed threshold is the fixed-sample instance of this rule \citep{nourbakhsh2026tracer}.

\item \textbf{Banking KPI-tuned evaluation protocol} (\S\ref{sec:evaluation}).
We wire outcome classes to production telephony telemetry---containment,
abandonment, handoff, compliance-escape---and define cost-weighted
detection-latency metrics that extend the early-detection score of \citet{nourbakhsh2026tracer}.
The protocol evaluates on real Spanish call logs with leave-one-campaign-out validation,
and accounts the decision latency of the trajectory monitor against the published
floors $\Lambda^*$ and the cessation floor \citep{garcia2026latencyfloor,garcia2026cessationfloor}.

\item \textbf{Intervention coupling} (\S\ref{sec:intervention}).
We define how grammar states trigger actionable interventions---proactive
clarification, escalation to a human agent, fraud-flag routing---inside a
production orchestration graph, with every intervention required to clear the
cessation floor \citep{garcia2026cessationfloor}.
This connects conversation-internal early warning to the escalation policies that
existing work decides from queue state alone \citep{queuecontrol2026}.
\end{enumerate}

The remainder of this paper is organised as follows.
Section~\ref{sec:background} situates the work across agent-trajectory benchmarks,
dialogue process mining, outcome prediction, and contact-center analytics.
Section~\ref{sec:grammar} develops the trajectory grammar and the drift statistic.
Section~\ref{sec:stopping} states and proves the Trajectory Stopping Theorem.
Section~\ref{sec:evaluation} defines the KPI-tuned evaluation protocol.
Section~\ref{sec:intervention} specifies intervention coupling.
Sections~\ref{sec:open_problems} and~\ref{sec:discussion} present open problems and limitations.
Section~\ref{sec:conclusion} concludes.
